\documentclass[11pt]{article}
\usepackage[margin=1in]{geometry}
\usepackage{amsmath,amssymb,amsthm}

\newcommand{\Smax}{S_{\max}}
\newcommand{\pistar}{\pi^\star}
\newcommand{\pihatstar}{\hat\pi^\star}
\newcommand{\Vstar}{V^\star}
\newcommand{\Vhatstar}{\hat V^\star}
\newcommand{\Vpihat}{V^{\pihatstar}}

\title{Why $\gamma$ Appears in the Propagation Term of the MDP Corollary}
\author{}
\date{}

\begin{document}
\maketitle

\section{Setup}

Consider two discounted-cost MDPs $M$ and $\hat M$ on the same state and
action spaces with discount factor $\gamma \in (0,1)$.  The Bellman equation
for the value function of a policy $\pi$ under model $M$ is
\begin{equation}\label{eq:bellman}
  V^\pi(s) = c_M(s, \pi(s)) + \gamma \sum_{s' \in \mathcal{S}} V^\pi(s')\, P_M(s' \mid s, \pi(s)).
\end{equation}

Let $\pihatstar$ denote the optimal policy of $\hat M$ and $\Vhatstar$ the
optimal value function of $\hat M$.  We want to bound $\Vpihat(s) - \Vhatstar(s)$.

\section{Derivation}

\paragraph{Step 1: Bellman equations.}
$\Vpihat$ satisfies the Bellman equation under the \emph{true} model $M$
with the \emph{approximate} policy $\pihatstar$:
\begin{equation}\label{eq:V-pihat}
  \Vpihat(s) = c_M\bigl(s, \pihatstar(s)\bigr)
  + \gamma \sum_{s'} \Vpihat(s')\, P_M\bigl(s' \mid s, \pihatstar(s)\bigr).
\end{equation}

$\Vhatstar$ satisfies the Bellman optimality equation under $\hat M$:
\begin{equation}\label{eq:V-hat}
  \Vhatstar(s) = \hat c\bigl(s, \pihatstar(s)\bigr)
  + \gamma \sum_{s'} \Vhatstar(s')\, \hat P\bigl(s' \mid s, \pihatstar(s)\bigr),
\end{equation}
where equality holds because $\pihatstar$ is optimal for $\hat M$.

\paragraph{Step 2: Take the difference.}
Subtracting~\eqref{eq:V-hat} from~\eqref{eq:V-pihat}:
\begin{equation}\label{eq:diff-raw}
  \Vpihat(s) - \Vhatstar(s)
  = \bigl[c_M(s, a) - \hat c(s, a)\bigr]
  + \gamma \sum_{s'} \Vpihat(s')\, P_M(s' \mid s, a)
  - \gamma \sum_{s'} \Vhatstar(s')\, \hat P(s' \mid s, a),
\end{equation}
where we write $a = \pihatstar(s)$ for brevity.

\paragraph{Step 3: Add and subtract $\gamma \sum_{s'} \Vhatstar(s') P_M(s' \mid s, a)$.}
This separates the future-value terms into a ``propagated error'' piece and
a ``transition mismatch'' piece:
\begin{align}
  \Vpihat(s) - \Vhatstar(s)
  &= \underbrace{\bigl[c_M(s,a) - \hat c(s,a)\bigr]}_{\text{cost mismatch}}
  \notag \\[6pt]
  &\quad + \gamma \sum_{s'}
    \underbrace{\bigl[\Vpihat(s') - \Vhatstar(s')\bigr]}_{\text{error at next state}}\,
    P_M(s' \mid s, a)
  \label{eq:decomposed} \\[6pt]
  &\quad + \gamma \biggl[
    \sum_{s'} \Vhatstar(s')\, P_M(s' \mid s, a)
    - \sum_{s'} \Vhatstar(s')\, \hat P(s' \mid s, a)
  \biggr].
  \notag
\end{align}

\paragraph{Step 4: Take absolute values and bound.}
Taking absolute values on both sides:
\begin{align}
  \bigl|\Vpihat(s) - \Vhatstar(s)\bigr|
  &\le \bigl|c_M(s,a) - \hat c(s,a)\bigr|
  \notag \\[4pt]
  &\quad + \gamma \sum_{s'} \bigl|\Vpihat(s') - \Vhatstar(s')\bigr|\, P_M(s' \mid s, a)
  \label{eq:bound} \\[4pt]
  &\quad + \gamma\, \biggl|
    \sum_{s'} \Vhatstar(s')\, P_M(s' \mid s, a)
    - \sum_{s'} \Vhatstar(s')\, \hat P(s' \mid s, a)
  \biggr|.
  \notag
\end{align}

\paragraph{Step 5: Identify the three terms.}
Define:
\begin{align*}
  \varepsilon(s, a) &= \bigl|c_M(s,a) - \hat c(s,a)\bigr|,
  \\[4pt]
  \Delta(s, a) &= \biggl|
    \sum_{s'} \Vhatstar(s')\, P_M(s' \mid s, a)
    - \sum_{s'} \Vhatstar(s')\, \hat P(s' \mid s, a)
  \biggr|.
\end{align*}
If $\alpha(s)$ is any function satisfying
$\bigl|\Vpihat(s) - \Vhatstar(s)\bigr| \le \alpha(s)$ for all $s$, then
\eqref{eq:bound} gives:
\begin{equation}\label{eq:beta-final}
  \boxed{%
    \alpha(s)
    \;\le\; \underbrace{\varepsilon(s,a)}_{\text{cost mismatch}}
    \;+\; \underbrace{\gamma \sum_{s'} \alpha(s')\, P_M(s' \mid s, a)}_{\text{propagated future error}}
    \;+\; \underbrace{\gamma\, \Delta(s,a)}_{\text{transition mismatch}}
  }
\end{equation}

The right-hand side defines the error certificate:
\begin{equation}\label{eq:beta-def}
  \beta(s,a)
  \;=\; \varepsilon(s,a)
  + \gamma \sum_{s'} \alpha(s')\, P_M(s' \mid s, a)
  + \gamma\,\Delta(s,a).
\end{equation}

\paragraph{Step 6: From $\beta$ to $\alpha$.}
The bound~\eqref{eq:bound} holds for any action $a$.  In particular, it holds
for $a = \pistar(s)$ (optimal in $M$) and $a = \pihatstar(s)$ (optimal in
$\hat M$).  The suboptimality gap $\Vpihat(s) - \Vstar(s)$ can be written as:
\begin{align*}
  \Vpihat(s) - \Vstar(s)
  &= \bigl[\Vpihat(s) - \Vhatstar(s)\bigr]
   + \bigl[\Vhatstar(s) - \Vstar(s)\bigr].
\end{align*}
The first bracket involves the policy $\pihatstar$ (the approximate
optimiser applied in the true model), and the second involves $\pistar$
(the true optimiser compared to the approximate value).  The proof of
Theorem~2 bounds these two brackets separately, each evaluated at its
respective action, giving:
\begin{equation}\label{eq:gap-two-actions}
  \Vpihat(s) - \Vstar(s)
  \;\le\; \beta\bigl(s,\, \pihatstar(s)\bigr) + \beta\bigl(s,\, \pistar(s)\bigr)
  \;\le\; 2\,\alpha(s),
\end{equation}
where
\begin{equation}\label{eq:alpha-max}
  \boxed{%
    \alpha_{\max}(s)
    \;=\; \max\!\Big\{
      \beta\bigl(s,\, \pistar(s)\bigr),\;\;
      \beta\bigl(s,\, \pihatstar(s)\bigr)
    \Big\}.
  }
\end{equation}
This is the \textbf{tighter variant}: $\alpha$ only needs to dominate $\beta$
at the two actions that actually appear in the gap decomposition, not at every
action in $\mathcal{A}$.

A \textbf{looser but simpler variant} maximises over all actions:
\begin{equation}\label{eq:alpha-sup}
  \alpha_{\sup}(s)
  \;=\; \max_{a \in \mathcal{A}}\, \beta(s, a)
  \;\ge\; \alpha_{\max}(s).
\end{equation}

Both define valid fixed-point equations.  The fixed-point iteration
initialises $\alpha^0(s) = 0$ and updates:
\[
  \alpha^{k+1}(s) = \max_{a \in A(s)}\, \beta^k(s,a),
\]
where $A(s) = \{\pistar(s),\, \pihatstar(s)\}$ for $\alpha_{\max}$
and $A(s) = \mathcal{A}$ for $\alpha_{\sup}$.

\section{Why $\gamma$ appears}

The $\gamma$ in front of $\sum_{s'} \alpha(s') P_M(s' \mid s, a)$ is
inherited directly from the Bellman equation~\eqref{eq:bellman}.  In a
discounted MDP, next-period values are multiplied by $\gamma$: an error of
magnitude $\alpha(s')$ one step in the future contributes only
$\gamma \cdot \alpha(s')$ to the current-period value.

This is also why the fixed-point operator
\[
  (\mathcal{T}\alpha)(s)
  = \max_a \bigl[
    \varepsilon(s,a) + \gamma \textstyle\sum_{s'} \alpha(s') P_M(s'|s,a) + \gamma\,\Delta(s,a)
  \bigr]
\]
is a $\gamma$-contraction in sup-norm: the only term that depends on $\alpha$
is multiplied by $\gamma < 1$.

\paragraph{Contrast with the finite-horizon theorem.}
In the general finite-horizon result (Theorem~2 of the paper), costs are
undiscounted and the recursion reads
\[
  \beta_t(h_t, a_t)
  = \varepsilon_t(h_t, a_t)
  + \int \alpha_{t+1}(h_{t+1})\, P(dh_{t+1} \mid h_t, a_t)
  + \Delta_t(h_t, a_t).
\]
No $\gamma$ appears because there is no discounting.  When specialising to
infinite-horizon discounted MDPs and taking the fixed-point limit ($T \to
\infty$), the discount factor $\gamma$ from the Bellman operator emerges
naturally in front of the propagation and transition-mismatch terms.

\end{document}
